\section{Surface Integrals and Flux}
\label{sec:surface-integrals}

A surface integral accumulates a scalar field over a two-dimensional surface or measures the flux of a vector field through that surface. Let
\[
\mathbf r(u,v)
\]
parametrize a surface. The tangent vectors are $\mathbf r_u$ and $\mathbf r_v$, and their cross product determines the oriented area element,
\[
\boxed{d\mathbf A=(\mathbf r_u\times\mathbf r_v)\,du\,dv.}
\]
Its magnitude gives the scalar area element,
\[
dA=\lVert\mathbf r_u\times\mathbf r_v\rVert\,du\,dv.
\]

For a scalar field $g$,
\[
\iint_S g\,dA
\]
weights each small surface patch by the local value of $g$. For a vector field $\mathbf F$, the flux is
\[
\boxed{\Phi_S=\iint_S\mathbf F\cdot d\mathbf A.}
\]
The dot product selects the component normal to the surface.

\begin{bookfigure}{Flux through an oriented surface. Only the normal component of the field contributes to $\mathbf F\cdot d\mathbf A$.}
\begin{tikzpicture}[scale=.9]
\draw[fill=visualcyan!8,draw=visualcyan!60] (-3,-.7) .. controls (-1,1.2) and (1,-1.2) .. (3,.7) -- (3,1.4) .. controls (1,-.5) and (-1,1.9) .. (-3,0) -- cycle;
\coordinate (P) at (.4,.45);
\draw[force vector] (P)--++(0,2) node[above] {$d\mathbf A$};
\draw[field vector] (P)--++(1.7,1.1) node[right] {$\mathbf F$};
\draw[dashed] (P)--++(0,1.15);
\node at (0,-1.3) {oriented surface $S$};
\end{tikzpicture}
\label{fig:surface-flux}
\end{bookfigure}

\subsection{Orientation}
A surface has two possible orientations. Reversing the normal changes $d\mathbf A$ to $-d\mathbf A$ and therefore reverses the sign of flux. For a closed surface the standard orientation is outward.

\begin{workedexample}
\textbf{Flux through a horizontal disk.}\par
Let $\mathbf F=(0,0,z)$ and let $S$ be the disk $x^2+y^2\leq a^2$ in the plane $z=h$, oriented upward. Since $\hat{\mathbf n}=\mathbf e_z$,
\[
\mathbf F\cdot\hat{\mathbf n}=h.
\]
Therefore
\[
\Phi_S=\iint_S h\,dA=h\pi a^2.
\]
The result is field strength times projected area because the field is uniform over the disk.
\end{workedexample}

\begin{mistakebox}[Using $dA$ when orientation matters]
The scalar element $dA$ has no direction. Flux requires the vector element $d\mathbf A=\hat{\mathbf n}dA$.
\end{mistakebox}
